% Figura 5: abertura real (em escala). C na origem, R = 6,5 cm.
% Abertura total phi = 10°: o espelho vai de -5° a 5°; do vértice ao topo, theta = 5°.
\begin{tikzpicture}[font=\small]
  \def\R{6.5}
  \draw[eixo] (-0.4,0) -- (7.6,0);
  \draw[line width=0.5pt,black!30,dashed] (5:\R) arc[start angle=5,end angle=12,radius=\R];
  \draw[line width=0.5pt,black!30,dashed] (-5:\R) arc[start angle=-5,end angle=-12,radius=\R];
  \draw[espelho] (-5:\R) arc[start angle=-5,end angle=5,radius=\R];
  \foreach \a in {-4.5,-3,...,4.5}
    \draw[hachura] (\a:\R) -- ++(\a-40:0.18);
  \draw[line width=0.6pt] (0,0) -- (5:\R);
  \draw[line width=0.4pt] (0:4.6) arc[start angle=0,end angle=5,radius=4.6];
  \node[anchor=base] at (3.0,0.42) {$R$};
  \node[inner sep=0pt,anchor=base west] at (4.75,0.1) {$\theta$};
  \draw[guia] (5:\R) -- (7.35,{\R*sin(5)});
  \draw[cota] (7.2,0) -- (7.2,{\R*sin(5)});
  \node[inner sep=0pt,anchor=base west] at (7.337,0.2) {$H$};
  \fill (0,0) circle (0.053);
  \node[anchor=base] at (-0.0,-0.435) {$C$};
  \node[inner sep=0pt,anchor=base west] at (6.656,-0.395) {$V$};
  \node[anchor=base] at (3.4,-0.988)
    {$\varphi \le 10^\circ \implies \theta = \varphi/2 \le 5^\circ \implies H = R\sin\theta \lesssim 0{,}09\,R$};
\end{tikzpicture}
